% s03_rows. Section file; new macros at the top; label prefix: row
\section{The closed form of the rows}
\label{sec:row}

\Cref{thm:row-pf,cor:row-unique,lem:row-closed12,cor:row-carry,cor:row-closed,prop:row-sympl} are statements about the closed form \eqref{eq:set-S} and \cref{def:set-rows} only, and they are proved by algebra on it, with no Gromov--Witten input. \Cref{prop:row-modrel} is a statement about the series \eqref{eq:set-gens} only. These results are all that \cref{sec:lea} uses from this section. The formal stationary phase calculus of \cref{sec:row-sp} and \cref{lem:row-crit,lem:row-analytic,lem:row-pfint,lem:row-wkb} are used only for the frame formula \cref{prop:fra-frame}, through \cref{prop:gen-osc,lem:spe-row0,lem:spe-red}. The identification of row $0$ of \cref{def:set-rows} with the canonical row of \cite{BS} is not made here: it is \cref{lem:spe-row0,lem:spe-red}.

%-------------------------------------------------------------------------------
\subsection{Picard--Fuchs for the closed form}
\label{sec:row-pf}

Fix a colour $i$ on $\TF$ and write $V=V_i=Ev_i$, $Y=Y_i$. Put $\mathsf{a}\coloneqq1+E/z$ and
\[
\mathsf{e}_n\coloneqq\frac{V^n}{(\mathsf{a})_n}\qquad(n\in\bbZ),
\]
with $(\mathsf{a})_n$ the Pochhammer symbol. Since $1/(\mathsf{a})_n=\prod_{j=1}^{n}z/(E+jz)$ for $n\geq 0$, we have
\begin{equation}
\label{eq:row-vS}
v\,\Srow(v,z)=z^{-1}\sum_{n\geq1}r_{n-1}\mathsf{e}_n .
\end{equation}
Let $\Pconj\coloneqq\re^{-\mathsf{u}_i/z}\LPF\re^{\mathsf{u}_i/z}$ at $\alpha=\alpha^*$, with $\theta\mathsf{u}_i=Y_i$; that is, $\Pconj$ is $\LPF$ with $z\theta$ replaced by $z\theta+Y_i$.

\begin{thm}\genusstatus{AI-proposed, unchecked}
\label{thm:row-pf}
Let $\mathsf{s}\mathsf{e}_n=\mathsf{e}_{n-1}$, $\mathsf{N}\mathsf{e}_n=n\mathsf{e}_n$ and $\mathsf{t}\coloneqq3\mathsf{s}-\mathsf{N}$. Then $\Pconj=(z^3/27)(\mathsf{t}^3-27\mathsf{s}^3)$, and
\begin{equation}
\label{eq:row-pfcoef}
\Pconj\Big(\sum_nb_n\mathsf{e}_n\Big)=-\frac{z^3}{27}\sum_m\big[27(m+1)b_{m+2}-3(3m^2+3m+1)b_{m+1}+m^3b_m\big]\mathsf{e}_m .
\end{equation}
Hence $\Delta_i^{-1/2}\,\hb(z)\,\Srow(v_i,z)\,\re^{\mathsf{u}_i/z}$ solves the Picard--Fuchs equation \eqref{eq:set-pf} at $\alpha=\alpha^*$ identically, for $i=0,1,2$.
\end{thm}

\begin{proof}
At every colour one has
\[
\theta V=1-\frac V3,\qquad Y+\frac E3=\frac EV,\qquad y(V-3)^3=1.
\]
Put $\mathsf{T}\coloneqq z\theta+Y+E/3$. Conjugation by $\re^{\mathsf{u}_i/z}$ gives
\[
\Pconj=\mathsf{T}^3+27y\prod_{j=0}^{2}\Big(\mathsf{T}-\frac E3+\frac{jz}{3}\Big).
\]
From $V\mathsf{e}_n=(\mathsf{a}+n)\mathsf{e}_{n+1}$ and $z(\mathsf{a}+n-1)=E+nz$ we get $\mathsf{T}\mathsf{e}_n=z\mathsf{e}_{n-1}-(nz/3)\mathsf{e}_n$. That is, $\mathsf{T}=(z/3)\mathsf{t}$, and
\[
\mathsf{T}-\frac E3+\frac{jz}3=-\frac13(\mathsf{O}-jz),\qquad\mathsf{O}\coloneqq E+z\mathsf{N}-3z\mathsf{s}.
\]
Let $\mathsf{M}$ be multiplication by $z(V-3)$. Then $\mathsf{O}=\mathsf{M}\mathsf{s}$ and $[\mathsf{s},\mathsf{M}]=z$, so $\mathsf{M}^2\mathsf{s}^2=\mathsf{O}(\mathsf{O}-z)$ and $\mathsf{M}^3\mathsf{s}^3=\mathsf{O}(\mathsf{O}-z)(\mathsf{O}-2z)$. Since $y\mathsf{M}^3=z^3$, this gives $\Pconj=(z^3/27)(\mathsf{t}^3-27\mathsf{s}^3)$.

Next,
\[
\mathsf{t}^3\mathsf{e}_n=27\mathsf{e}_{n-3}-27(n-1)\mathsf{e}_{n-2}+3(3n^2-3n+1)\mathsf{e}_{n-1}-n^3\mathsf{e}_n,
\]
which gives the coefficient in \eqref{eq:row-pfcoef}. With $b_n=r_{n-1}$ this coefficient vanishes by \eqref{eq:set-rec} for $m\geq0$, and trivially for $m\leq-1$.

\emph{Convergence, $z$-adic form.} Since $1/(\mathsf{a})_n=\prod_{j\leq n}z/(E+jz)$, each $\mathsf{e}_n$ with $n\geq0$ lies in $V^nz^n\bbQ(E)\fps{z}$, so $\sum b_n\mathsf{e}_n$ converges $z$-adically. The operator $\Pconj$ has coefficients polynomial in $z$ and never lowers the $z$-order. So $[z^k]\Pconj(v\Srow)$ is a finite expression, a Laurent polynomial in $V$ and $(V-3)$, and it vanishes in $\bbQ(\epsilon_1,\epsilon_2)(t_i)$. This is the form used in \cref{prop:row-sympl}.

\emph{Convergence, $V$-adic form, at the conifold.} The same computation holds coefficientwise in $V$, with coefficients rational in $z$. On the $V$-adic side, $\mathsf{e}_n$ lies in $V^n\bbQ(\epsilon_1,\epsilon_2)(z)$, the image $\Pconj\mathsf{e}_n$ is the finite combination of $\mathsf{e}_{n-2},\mathsf{e}_{n-1},\mathsf{e}_n$ given by \eqref{eq:row-pfcoef}, and $y=(V-3)^{-3}$ is analytic at $V=0$. So $\Pconj(v\Srow)=0$ holds in $\bbQ(\epsilon_1,\epsilon_2)(z)\fps{V}$. This is the form used in \eqref{eq:row-zinf}.

Finally, $\Delta_i^{-1/2}=v_i/\sqrt{3\epsilon_1\epsilon_2E}$, so $\Delta_i^{-1/2}\Srow$ is a constant times $v\Srow=z^{-1}\sum b_n\mathsf{e}_n$, and the factors $\hb(z)$, $z^{-1}$, which depend on $z$ only, commute with $\Pconj$.
\end{proof}

\begin{cor}\genusstatus{AI-proposed, unchecked}
\label{cor:row-unique}
If $b_n\in\bbC\lser{z}$, $b_n=0$ for $n\leq0$, and $\Pconj(\sum b_n\mathsf{e}_n)=0$, the sum being understood $V$-adically, then $b_n=b_1(z)r_{n-1}$, so that $\sum b_n\mathsf{e}_n=z\,b_1(z)\,v\Srow$. Thus the Picard--Fuchs equation fixes the row only up to a factor depending on $z$ alone. If the $b_n$ are constants and $\Pconj(\sum b_n\mathsf{e}_n)=0$ holds at every $z$-order (each order being a rational function of $V$), the same conclusion follows.
\end{cor}

\begin{proof}
By \cref{thm:row-pf} the expression is $-(z^3/27)\sum_mC_m\mathsf{e}_m$ with
\[
C_m\coloneqq27(m+1)b_{m+2}-3(3m^2+3m+1)b_{m+1}+m^3b_m.
\]
The terms with $m<0$ vanish, since $b_n=0$ for $n\leq0$. Further $\mathsf{e}_m=(Vz/E)^m/\prod_{j\leq m}(1+jz/E)$, so the coefficient of $V^m$ is $-(z^3/27)\,C_m\,(z/E)^m/\prod_{j\leq m}(1+jz/E)$. Hence $C_m=0$ for all $m\geq0$. This is the recursion \eqref{eq:set-rec} for $b_{n}/b_1$, with $b_0=0$. In the second case, $\mathsf{e}_m=(Vz/E)^m(1+O(z))$, so if $C_0=\cdots=C_{m-1}=0$ the lowest $z$-order of $\sum_kC_k\mathsf{e}_k$ has coefficient $C_m(V/E)^m$; hence $C_m=0$ by induction on $m$, and the recursion follows as before.
\end{proof}

%-------------------------------------------------------------------------------
\subsection{Formal stationary phase}
\label{sec:row-sp}

The contents of this subsection are formal and hold in $n$ dimensions. They are not used in \cref{sec:lea,sec:van,sec:val} or in the regularity theorems.

\emph{The calculus.} Let a phase $\phi$ have a non-degenerate critical point, placed at $\zeta=0$ in coordinates $\zeta=(\zeta^1,\ldots,\zeta^n)$, with Hessian $\Hess_{ab}$ and critical value $\phi_0$. Write $\phi=\phi_0+\frac12\Hess_{ab}\zeta^a\zeta^b+\phi^{\mathrm{ng}}$. Let $\mathscr{R}$ be the ring of sums $\sum c_{k,\mathbf{b}}z^k\zeta^{\mathbf{b}}$ (with $\mathbf{b}$ a multi-index) having, for each $N$, finitely many terms with $k+|\mathbf{b}|/2\leq N$. Put
\[
\langle f\rangle\coloneqq\re^{\Delta_{\Hess}}f\big|_{\zeta=0},\qquad\Delta_{\Hess}\coloneqq\frac z2\Hess^{ab}\partial_a\partial_b.
\]
Since $\Delta_{\Hess}$ preserves $k+|\mathbf{b}|/2$, this is continuous on $\mathscr{R}$. For an amplitude $\gamma$ analytic at $\zeta=0$, define
\[
\Lambda(\gamma)\coloneqq(\det\Hess)^{-1/2}\big\langle\re^{-\phi^{\mathrm{ng}}/z}\gamma\big\rangle.
\]
This is the normalised expansion \eqref{eq:set-asym}: $\Asym[\re^{-\phi/z}\gamma\,\dd^n\zeta]=\Lambda(\gamma)$ at the critical point. The calculus is written for $\re^{-\phi/z}$. The integrand $\re^{\cW/z}$ corresponds to $\phi=-\cW$. In dimension one, \cref{lem:row-SP5} with $z\mapsto-z$ is how the factors $\Theta(-z/a)$ arise. The difference of Gaussian conventions is the constant $\mathsf{k}$ of \cref{prop:gen-osc}.

\begin{lem}[Gaussian integration by parts]\genusstatus{AI-proposed, unchecked}
\label{lem:row-G1}
$\langle(\Hess\zeta)_af\rangle=z\langle\partial_af\rangle$ for $f\in\mathscr{R}$.
\end{lem}

\begin{proof}
The commutator $[\Delta_{\Hess},\zeta^b]=z\Hess^{bc}\partial_c$ commutes with $\Delta_{\Hess}$. So $\re^{\Delta_{\Hess}}\zeta^b=(\zeta^b+z\Hess^{bc}\partial_c)\re^{\Delta_{\Hess}}$, and at $\zeta=0$ this gives $\langle\zeta^bf\rangle=z\Hess^{bc}\langle\partial_cf\rangle$. Contract with $\Hess_{ab}$.
\end{proof}

\begin{lem}[Total derivatives vanish]\genusstatus{AI-proposed, unchecked}
\label{lem:row-SP2}
$\Lambda\big(\partial_aB-z^{-1}(\partial_a\phi)B\big)=0$ for every analytic $B$.
\end{lem}

\begin{proof}
Let $f\coloneqq\re^{-\phi^{\mathrm{ng}}/z}B$. Since $\partial_a\phi=(\Hess\zeta)_a+\partial_a\phi^{\mathrm{ng}}$, we have $\re^{-\phi^{\mathrm{ng}}/z}(\partial_aB-z^{-1}(\partial_a\phi)B)=\partial_af-z^{-1}(\Hess\zeta)_af$. Apply \cref{lem:row-G1}.
\end{proof}

\begin{lem}[Fixed critical point]\genusstatus{AI-proposed, unchecked}
\label{lem:row-FP}
Let $\phi_p,\gamma_p$ depend on a parameter $p$. If the critical point sits at $\zeta=0$ for every $p$ (the linear term of $\phi_p$ vanishes identically in $p$) and the critical value does not depend on $p$, then
\[
\frac{\dd}{\dd p}\Lambda_p(\gamma_p)=\Lambda_p\big(\gamma_p'-z^{-1}\phi_p'\gamma_p\big).
\]
\end{lem}

\begin{proof}
Differentiation is termwise legitimate, because each $z$-coefficient is a finite sum, by the formula in the proof of \cref{lem:row-analytic}. Let $f\coloneqq\re^{-\phi^{\mathrm{ng}}/z}\gamma$. With $\Delta_{\Hess}'\coloneqq\dd\Delta_{\Hess}/\dd p=-\frac z2(\Hess^{-1}\Hess'\Hess^{-1})^{ab}\partial_a\partial_b$, which commutes with $\Delta_{\Hess}$, we have
\[
\frac{\dd}{\dd p}\big[(\det\Hess)^{-1/2}\langle f\rangle\big]=(\det\Hess)^{-1/2}\Big[-\tfrac12\operatorname{tr}(\Hess^{-1}\Hess')\langle f\rangle+\langle\Delta_{\Hess}'f\rangle+\langle f'\rangle\Big].
\]
Since $\phi_p'=(\phi^{\mathrm{ng}})'+\frac12\Hess_{ab}'\zeta^a\zeta^b$, the claim is equivalent to
\[
-\frac{1}{2z}\Hess_{ab}'\langle\zeta^a\zeta^bf\rangle=-\tfrac12\operatorname{tr}(\Hess^{-1}\Hess')\langle f\rangle+\langle\Delta_{\Hess}'f\rangle.
\]
Applying \cref{lem:row-G1} twice gives $\langle\zeta^a\zeta^bf\rangle=z\Hess^{ab}\langle f\rangle+z^2\Hess^{ac}\Hess^{bd}\langle\partial_c\partial_df\rangle$. Substituting reproduces the two terms $-\frac12\operatorname{tr}(\Hess^{-1}\Hess')$ and $\Delta_{\Hess}'$.
\end{proof}

\begin{lem}[Moving critical point]\genusstatus{AI-proposed, unchecked}
\label{lem:row-SP3}
Let $x_0(p)$ be a non-degenerate critical point of $\phi_p$ depending analytically on $p$, with critical value $\phi_0(p)$, and let $\gamma_p$ be any analytic amplitude. Then
\[
\Big(\frac{\dd}{\dd p}-z^{-1}\frac{\dd\phi_0}{\dd p}\Big)\Lambda_p(\gamma_p)=\Lambda_p\big(\partial_p\gamma-z^{-1}(\partial_p\phi)\gamma\big),
\]
that is, $\dd_p[\re^{-\phi_0/z}\Lambda_p(\gamma_p)]=\re^{-\phi_0/z}\Lambda_p(\partial_p\gamma-z^{-1}(\partial_p\phi)\gamma)$. Here $\partial_p$ is taken at fixed $x$.
\end{lem}

\begin{proof}
Translate the critical point to $0$: put $\zeta\coloneqq x-x_0(p)$, $\widetilde\phi_p(\zeta)\coloneqq\phi_p(x_0(p)+\zeta)-\phi_0(p)$ and $\widetilde\gamma_p(\zeta)\coloneqq\gamma_p(x_0(p)+\zeta)$. \Cref{lem:row-FP} applies to these. We have $\widetilde\gamma'=\partial_p\gamma+x_0'^a\partial_a\gamma$ and $\widetilde\phi'=\partial_p\phi+x_0'^a\partial_a\phi-\dd\phi_0/\dd p$. So the right side of \cref{lem:row-FP} is
\[
\Lambda_p\big(\partial_p\gamma-z^{-1}(\partial_p\phi)\gamma\big)+z^{-1}\frac{\dd\phi_0}{\dd p}\Lambda_p(\gamma)+x_0'^a\,\Lambda_p\big(\partial_a\gamma-z^{-1}(\partial_a\phi)\gamma\big),
\]
and the terms in $x_0'$ vanish by \cref{lem:row-SP2}.
\end{proof}


\begin{lem}[Coordinate invariance]\genusstatus{AI-proposed, unchecked}
\label{lem:row-SP1}
Let $x=g(\zeta)$ be a local biholomorphism with $g(0)=x_0$ a non-degenerate critical point of $\phi$. With the branch convention $(\det\widetilde{\Hess})^{1/2}=\det\dd g(0)\cdot(\det\Hess)^{1/2}$, where $\widetilde{\Hess}$ is the Hessian of $\phi\circ g$ at $0$, the expansion of $\phi\circ g$ with amplitude $(\gamma\circ g)\det\dd g$ equals the expansion of $\phi$ with amplitude $\gamma$.
\end{lem}

\begin{proof}
The linear case is direct. In general, interpolate by $g_s(\zeta)\coloneqq x_0+s^{-1}(g(s\zeta)-x_0)$, with $g_0(\zeta)\coloneqq x_0+\dd g(0)\zeta$. Put $\phi_s\coloneqq\phi\circ g_s$ and $\gamma_s\coloneqq(\gamma\circ g_s)\det\dd g_s$. Since $\dd g_s(0)=\dd g(0)$, the Hessian $\Hess_s$ is constant in $s$. Since $g_s(0)=x_0$ for all $s$, the vector field $X_s\coloneqq(\dd g_s)^{-1}\partial_sg_s$ vanishes at $0$. We have $\phi_s'=X_s(\phi_s)$ and $\gamma_s'=X_s(\gamma_s)+\operatorname{div}(X_s)\gamma_s$. Let $\mathcal{I}(s)$ be the expansion at $s$. By \cref{lem:row-FP},
\[
\frac{\dd}{\dd s}\mathcal{I}(s)=\Lambda_s\big(X_s(\gamma_s)+\operatorname{div}(X_s)\gamma_s-z^{-1}X_s(\phi_s)\gamma_s\big)=\sum_a\Lambda_s\big(\partial_a(X_s^a\gamma_s)-z^{-1}(\partial_a\phi_s)X_s^a\gamma_s\big),
\]
which vanishes by \cref{lem:row-SP2}. Each $z$-coefficient of $\mathcal{I}$ is polynomial in $s$, so $\mathcal{I}(1)=\mathcal{I}(0)$, and the latter is the linear case.
\end{proof}

\begin{lem}[Fubini]\genusstatus{AI-proposed, unchecked}
\label{lem:row-SP4}
If the phase splits exactly, $\phi(x',x'')=\phi'(x')+\phi''(x'')$, then the operator $\Delta_{\Hess}$ and the non-Gaussian exponential split into commuting parts, and $\Lambda$ is the iterated expansion in $x'$ and $x''$. The amplitude need not split.
\end{lem}

\begin{proof}
The Hessian is block diagonal, so $\Delta_{\Hess}=\Delta_{\Hess'}+\Delta_{\Hess''}$ with commuting summands, $\det\Hess=\det\Hess'\det\Hess''$ and $\phi^{\mathrm{ng}}=\phi'^{\mathrm{ng}}+\phi''^{\mathrm{ng}}$. Each operator acts on the amplitude regarded as a function of one group of variables, with the other group as parameters.
\end{proof}

\begin{lem}[Formal Stirling]\genusstatus{AI-proposed, unchecked}
\label{lem:row-SP5}
In this lemma $a$ denotes a non-zero number and $\beta$ a number. Then
\[
\Asym_{q=a}\big[\re^{-(q-a\log q)/z}\,q^{\beta-1}\,\dd q\big]=a^{\beta-1/2}\,\Theta(z/a,\beta),
\]
with the branch of $a^{\beta-1/2}$ fixed by the convention of \cref{lem:row-SP1}; this matters at $a=-E$.
\end{lem}

\begin{proof}
Scaling $q=ap$ (\cref{lem:row-SP1}) gives the form $a^{\beta-1/2}\mathcal{S}(w,\beta)$ with $w=z/a$. \Cref{lem:row-SP2} with $B=q^\beta$ gives
\begin{equation}
\tag{S}
\mathcal{S}(w,\beta+1)=(1+\beta w)\mathcal{S}(w,\beta).
\end{equation}
\Cref{lem:row-SP3} with $p=a$ gives $z\partial_a\mathcal{I}_a+(\log a)\mathcal{I}_a=\partial_\beta \mathcal{I}_a$ for $\mathcal{I}_a\coloneqq a^{\beta-1/2}\mathcal{S}$, since $\partial_a\phi=-\log q$, $\dd\phi_0/\dd a=-\log a$, and $\log q$ under $\Lambda$ is $\partial_\beta$. The term $\log a$ cancels, and this reads
\begin{equation}
\tag{D}
\partial_\beta\mathcal{S}=(\beta-\tfrac12)w\mathcal{S}-w^2\partial_w\mathcal{S}.
\end{equation}
Each $z$-coefficient of $\mathcal{S}$ is a polynomial in $\beta$, because $(1+\zeta/a)^{\beta-1}$ has jets polynomial in $\beta$; this is what makes $\partial_\beta$ termwise legitimate in (D). Write $\log\mathcal{S}=\sum_{k\geq1}\ell_k(\beta)w^k$. The difference $d_k$ of two polynomial solutions has period $1$ by (S), so it is constant. Equation (D) reads $\partial_\beta\ell_1=\beta-\frac12$ and $\partial_\beta\ell_{k+1}=-k\ell_k$, which forces $d_k=0$ for all $k\geq1$. For $\beta$ not an integer, $a^\beta$ needs a chosen $\log a$; the applications have $\beta\in\bbZ$. The series $\Theta(w,\beta)$ satisfies (S) and (D), using $B_{k+1}'=(k+1)B_k$.
\end{proof}

%-------------------------------------------------------------------------------
\subsection{Critical points, analyticity, Picard--Fuchs under the integral}
\label{sec:row-crit}

In log coordinates $x_j=\log q_j$, the superpotential \eqref{eq:set-spot} is
\[
\cW=\sum_j\big(\re^{x_j}-\alpha_jx_j\big)+y^{-1/3}\re^{(x_0+x_1+x_2)/3},\qquad\dd\Omega=\tfrac13\dd^3x,
\]
and we put $\Pcr\coloneqq y^{-1/3}\re^{(x_0+x_1+x_2)/3}$.

\begin{lem}\genusstatus{AI-proposed, unchecked}
\label{lem:row-crit}
\begin{enumerate}
\item The critical points of $\cW$ are exactly the points with $q_j=\alpha_j+Y$, $\Pcr=-3Y$, where $\Lhat(Y)=0$ and $Y(\alpha_j+Y)\neq0$.
\item $\det\mathrm{Hess}\,\cW=W(Y)/3$, and on $\Lhat=0$ one has $W=-Y\Lhat'(Y)$.
\item At $\TF$ the three roots $Y_i$ are finite, non-zero, simple and non-degenerate if and only if $y\notin\{0,-1/27\}$.
\item At $y=-1/27$ only the colour-$0$ root escapes on the real branch, since $1+3\omega^it\to1-\omega^i\neq0$ for $i=1,2$.
\item $\theta\mathsf{u}=Y$ at each critical point, where $\mathsf{u}$ is the critical value.
\end{enumerate}
\end{lem}

\begin{proof}
We have $\partial_j\cW=q_j-\alpha_j+\Pcr/3$. Put $Y\coloneqq-\Pcr/3$. The constraint $\Pcr^3=\prod_jq_j/y$ becomes $\Lhat(Y)=0$. Lifts to log coordinates differ by $2\pi\ri\bbZ$ and have identical jets. The Hessian is $\Hess=\mathrm{diag}(q)-(Y/3)\mathsf{J}$, with $\mathsf{J}$ the all-ones matrix, and with $e_2(q)=e_2(\alpha)+2e_1(\alpha)Y+3Y^2$ this gives
\[
\det\Hess=\prod_jq_j-\frac Y3e_2(q)=\frac W3.
\]
From $81yY^2=-3\prod_j(Y+\alpha_j)/Y$ we get $Y\Lhat'=Ye_2(Y+\alpha)-3\prod_j(Y+\alpha_j)=-W$. On $\TF$, $(Y_i+E/3)/Y_i=-3t_i$ and $W(Y_i)=Eq_i^2$ by \eqref{eq:set-WTF}, which gives (iii) and (iv). Finally $\theta\mathsf{u}=\theta\cW|_{\mathrm{cr}}=-\Pcr/3=Y$.
\end{proof}

The algebraic row is the following series. Let $\Hess\coloneqq\mathrm{diag}(\alpha_j+Y)-(Y/3)\mathsf{J}$, so that $\det\Hess=W/3$, and let
\[
\Jet_m(\zeta)\coloneqq\sum_j(\alpha_j+Y)\frac{\zeta_j^m}{m!}-3Y\frac{(\sum_j\zeta_j/3)^m}{m!}.
\]
Then
\begin{equation}
\label{eq:row-Phi}
\Phi(\alpha,Y;z)\coloneqq\exp\Big(-\frac z2(\Hess^{-1})^{ab}\partial_a\partial_b\Big)\exp\Big(z^{-1}\sum_{m\geq3}\Jet_m(\zeta)\Big)\Big|_{\zeta=0}.
\end{equation}

\begin{lem}\genusstatus{AI-proposed, unchecked}
\label{lem:row-analytic}
On a parameter domain where the critical point $x_c(p)$ is non-degenerate and analytic, and with an analytic branch of $(\det\Hess)^{-1/2}$ on a simply connected set, every $z$-coefficient of $\Asym$ is analytic in $p$. For $\cW$: let $U$ be a simply connected open set of $(y,\alpha)$, $y\neq0$, on which a root $Y$ of $\Lhat$ is analytic with $YW(Y)\prod_j(Y+\alpha_j)\neq0$. On $U$,
\[
\Asy\coloneqq\Asym_{x^{\mathrm{cr}}}\big[\re^{\cW/z}\dd\Omega\big]=\mathsf{k}_0\,W^{-1/2}\,\Phi(\alpha,Y;z),
\]
with $\mathsf{k}_0^2=-1/3$, $[z^n]\Phi\in\bbQ[\alpha,Y,W^{-1}]$ and $[z^0]\Phi=1$.
\end{lem}

\begin{proof}
Write $\phi=-\cW$, with $\Hess^{\phi}$ its Hessian, and $\gamma$ for the amplitude. The coefficient of $z^n$ in $\re^{\Delta_{\Hess^\phi}}(\re^{-\phi^{\mathrm{ng}}/z}\gamma)|_{x_c}$ is the finite sum, over $m-k=n$ with $3k\leq2m$, of
\[
\big(2^m\,m!\,k!\big)^{-1}\big((\Hess^{\phi})^{ab}\partial_a\partial_b\big)^m\big[(-\phi^{\mathrm{ng}})^k\gamma\big]\Big|_{x_c}.
\]
The bound $3k\leq2m$ holds because $\phi^{\mathrm{ng}}$ vanishes to order $3$ at $x_c$; it gives $k\leq2n$. The sum is polynomial in the jets of $\phi$, in $\gamma$, and in $\mathrm{adj}(\Hess^\phi)/\det\Hess^\phi$. This gives the first assertion. For $\cW$, the derivatives $\partial^{\mathbf{b}}\cW$ at the critical point, with $|\mathbf{b}|=r\geq2$, are $q_j\delta_{\mathbf{b}=r\mathbf{1}_j}-Y/3^{r-1}$, with $\mathbf{1}_j$ the $j$-th unit multi-index; so the terms of order $m$ of the Taylor expansion form $\Jet_m(\zeta)$ for $m\geq2$, and $\Jet_2$ gives the Hessian $\Hess$ of $\cW$. Since $\Hess^\phi=-\Hess$ and $-\phi^{\mathrm{ng}}=\sum_{m\geq3}\Jet_m$, the finite sum above, divided by its $z^0$ term, is the series $\Phi$ of \eqref{eq:row-Phi}. As $\Hess^{-1}=(3/W)\,\mathrm{adj}\,\Hess$, this gives $[z^n]\Phi\in\bbQ[\alpha,Y,W^{-1}]$ and $[z^0]\Phi=1$. The constant $\mathsf{k}_0$ comes from $\dd\Omega=\frac13\dd^3x$, $\det\mathrm{Hess}(-\cW)=-W/3$ and $\det\mathrm{Hess}\,\cW=W/3$ (\cref{lem:row-crit}).
\end{proof}

In \cref{lem:row-pfint} the expression $\LPF(\re^{\mathsf{u}/z}f)$ stands for $\re^{\mathsf{u}/z}\Pconj f$, which needs only $\theta\mathsf{u}=Y$.
\begin{lem}\genusstatus{AI-proposed, unchecked}
\label{lem:row-pfint}
For every $(y,\alpha)$ in $U$, including $\alpha=\alpha^*$, the expression $\re^{\mathsf{u}/z}\Asy$ is annihilated by $\LPF$.
\end{lem}

\begin{proof}
For an amplitude $h(x,y)$ analytic near the critical point put
\[
\mathsf{I}[h]\coloneqq\re^{\mathsf{u}/z}\Asym\big[h\,\re^{\cW/z}\dd\Omega\big].
\]
All amplitudes below are polynomials in the $q_j$ and $\Pcr$. Two rules follow from \cref{lem:row-SP3,lem:row-SP2}. Since $\theta\cW=-\Pcr/3$ at fixed $x$ and $\partial_j\cW=q_j-\alpha_j+\Pcr/3$, we have
\begin{align}
z\theta\,\mathsf{I}[h]&=\mathsf{I}\big[z\theta h-\tfrac13\Pcr h\big], \tag{A}\\
0&=\mathsf{I}\big[z\partial_jh+(q_j-\alpha_j+\tfrac13\Pcr)h\big]. \tag{B}
\end{align}
Rule (A) is \cref{lem:row-SP3} with $p=\log y$, for any analytic amplitude, and rule (B) is \cref{lem:row-SP2}. Adding (A) to (B) and to $\alpha_j\mathsf{I}[h]$ gives, for each $j$,
\begin{equation}
\tag{C}
(z\theta+\alpha_j)\mathsf{I}[h]=\mathsf{I}\big[q_jh+z\partial_jh+z\theta h\big].
\end{equation}
Here $\theta$ acts on the explicit $y$-dependence of $h$ at fixed $x$, and it commutes with $\partial_x$.

\emph{Step 1: from $\prod_j(z\theta+\alpha_j)$ to $q_0q_1q_2$.} Apply (C) three times, with amplitudes $1$, $q_2$, $q_1q_2$. Since $\partial_1q_2=0$, $\partial_0(q_1q_2)=0$ and $\theta q_j=0$ at fixed $x$,
\[
(z\theta+\alpha_2)\mathsf{I}[1]=\mathsf{I}[q_2],\qquad(z\theta+\alpha_1)\mathsf{I}[q_2]=\mathsf{I}[q_1q_2],\qquad(z\theta+\alpha_0)\mathsf{I}[q_1q_2]=\mathsf{I}[q_0q_1q_2].
\]
So $\prod_j(z\theta+\alpha_j)\mathsf{I}[1]=\mathsf{I}[q_0q_1q_2]$. This uses \cref{lem:row-SP3} three times and \cref{lem:row-SP2} three times, with no division.

\emph{Step 2: from $q_0q_1q_2$ to the third-order operator.} At fixed $x$, $\theta\Pcr^m=-(m/3)\Pcr^m$. By (A) with $h=\Pcr^m$ we get $z\theta\mathsf{I}[\Pcr^m]=\mathsf{I}[-(mz/3)\Pcr^m-\frac13\Pcr^{m+1}]$, that is,
\[
\mathsf{I}[\Pcr^{m+1}]=-z(3\theta+m)\mathsf{I}[\Pcr^m].
\]
Hence $\mathsf{I}[\Pcr^3]=-z^3(3\theta+2)(3\theta+1)(3\theta)\mathsf{I}[1]$, where the factors commute. This uses \cref{lem:row-SP3} three times.

\emph{Step 3.} We have $q_0q_1q_2=y\Pcr^3$, and multiplication by the scalar $y$ passes through $\mathsf{I}$ by linearity. So
\[
\prod_j(z\theta+\alpha_j)\mathsf{I}[1]=y\mathsf{I}[\Pcr^3]=-z^3y\prod_{j=0}^{2}(3\theta+j)\mathsf{I}[1],
\]
which is $\LPF\mathsf{I}[1]=0$.
\end{proof}

%-------------------------------------------------------------------------------
\subsection{Uniqueness}
\label{sec:row-wkb}

For generic or diagonal weights, let $\Pconj\coloneqq\re^{-\mathsf{u}/z}\LPF\re^{\mathsf{u}/z}=\sum_{m=0}^{3}z^m\Pconj_m$, with $\theta\mathsf{u}=Y$.

\begin{lem}\genusstatus{AI-proposed, unchecked}
\label{lem:row-wkb}
Let $\Omega'$ be a connected domain on which $Y$ is holomorphic, $\Lhat(Y)=0$ and $YW(Y)\neq0$. Then:
\begin{enumerate}
\item $\Pconj_0=0$ and $\Pconj_1=-(W/Y)\big(\theta+\frac12\theta\log W\big)$. Hence $\ker\Pconj_1=\bbC\cdot W^{-1/2}$.
\item If $f=\sum f_kz^k$ and $\mathsf{f}'=\sum\mathsf{f}'_kz^k$ have coefficients holomorphic on $\Omega'$, $\Pconj f=\Pconj\mathsf{f}'=0$ and $\mathsf{f}'_0$ is not identically $0$, then $f=\mathsf{c}(z)\mathsf{f}'$ for a unique $\mathsf{c}\in\bbC\fps{z}$. When $\alpha$ varies, $\mathsf{c}$ depends on $\alpha$ only, not on $y$.
\item If moreover $f,\mathsf{f}'$ extend analytically to a boundary point $v_{\flat}$ of $\Omega'$, with $\mathsf{f}'(v_{\flat})$ invertible in $\bbC\fps{z}$, then $\mathsf{c}=f(v_{\flat})/\mathsf{f}'(v_{\flat})$.
\end{enumerate}
\end{lem}

\begin{proof}
(i) Put $\mathsf{T}'\coloneqq z\theta+Y$. Then
\[
\Pconj=\prod_j(\mathsf{T}'+\alpha_j)+27y\mathsf{T}'^3+y\big(27z\mathsf{T}'^2+6z^2\mathsf{T}'\big),
\]
because $3\mathsf{T}'(3\mathsf{T}'+z)(3\mathsf{T}'+2z)=27\mathsf{T}'^3+27z\mathsf{T}'^2+6z^2\mathsf{T}'$. Since
\[
\mathsf{T}'^n\mathsf{f}=Y^n\mathsf{f}+z\Big[nY^{n-1}\theta\mathsf{f}+\binom n2Y^{n-2}(\theta Y)\mathsf{f}\Big]+O(z^2),
\]
we get $\Pconj_0=\Lhat(Y)=0$ and $\Pconj_1\mathsf{f}=\Lhat'(Y)\theta\mathsf{f}+\big[\frac12\Lhat''(Y)\theta Y+27yY^2\big]\mathsf{f}$, where $'$ is the $Y$-derivative at fixed $y$. Now $\Lhat'=-W/Y$ and $\theta Y=-Y\prod_j(Y+\alpha_j)/W=27yY^4/W$. Differentiating $W=-Y\Lhat'(Y)$ along $\theta$ gives $\theta W=(W/Y)\theta Y-Y\Lhat''\theta Y-81yY^3$. These give $\frac12\Lhat''\theta Y+27yY^2=-\theta W/(2Y)$, which is (i).

(ii) The coefficient of $z^{k+1}$ in $\Pconj f=0$ is $\Pconj_1f_k+\Pconj_2f_{k-1}+\Pconj_3f_{k-2}=0$, and $\Pconj$ commutes with series in $z$ alone. By (i), $f_0$ and $\mathsf{f}'_0$ are both multiples of $W^{-1/2}$, and the multiple for $\mathsf{f}'_0$ is non-zero. So $f_0=\mathsf{c}_0\mathsf{f}'_0$ for a constant $\mathsf{c}_0$, and $z^{-1}(f-\mathsf{c}_0\mathsf{f}')$ is again a solution with holomorphic coefficients. Induct.

(iii) Evaluate at $v_\flat$.
\end{proof}

The proof is algebraic. Parts (i) and (ii) hold verbatim with holomorphic functions on $\Omega'$ replaced by $\bbF\fps{y}$, with $\bbF$ a field of characteristic $0$, and $Y\in\bbF\fps{y}$ with $\Lhat(Y)=0$ and $Y(0)W(Y)(0)\neq0$. The $\theta$-constants are then $\bbF$, and $f,\mathsf{f}'$ are series in $z$ with coefficients in $\bbF\fps{y}$. In (i), $W^{-1/2}$ need not lie in $\bbF\fps{y}$. The statement there is that every non-zero element of $\ker\Pconj_1$ on $\bbF\fps{y}$ is a unit (compare $y$-orders, using $\theta\log W\in y\bbF\fps{y}$) and that two such elements have ratio in $\bbF$, which is all that (ii) uses. With $\bbF=\Kalh$ of \cref{prop:gen-osc}, $\ker\Pconj_1=\bbF\,(3W)^{-1/2}$.

%-------------------------------------------------------------------------------
\subsection{Rows 1 and 2 in closed form}
\label{sec:row-closed12}

On $\TF$, write $v=v_i$, $Y=Y_i$ at a fixed colour $i$. Then
\begin{equation}
\label{eq:row-thetav}
\theta v=\frac1E-\frac v3,\qquad Y=-\frac E3+\frac1v,\qquad\theta\log v=\frac YE,\qquad\Delta_i\propto v^{-2}.
\end{equation}

\begin{lem}\genusstatus{AI-proposed, unchecked}
\label{lem:row-closed12}
For any non-vanishing section $\cone$,
\begin{align*}
\Ph_1&=E\,\cone^{-1}\,\frac{\theta v}{v}\,\Srow^{[1]},\\
\Ph_2&=E\,\cone\,\xi\,\theta v\Big[\frac{E\,\theta v\,\Srow^{[2]}}{v^2}+\big(z\,\theta\log(\cone^{-1}\theta v)-\dtwo\big)\frac{\Srow^{[1]}}{v}\Big].
\end{align*}
\end{lem}

\begin{proof}
From \eqref{eq:row-thetav} we have the identity $Y+z\theta\log v=(E+z)\theta v/v$. Conjugating the operators of \cref{def:set-rows} by $\Delta_i^{1/2}\propto v^{-1}$ replaces $\Thp_i$ by $\Thp_i+z\theta\log v$ on the normalised rows. On the $l$-th term $r_l(vz)^l/\prod_{j\leq l+1}(1+jz/E)$ of $\Srow$ this operator acts as multiplication by $(E+(l+1)z)\theta v/v$, and
\[
\frac{E+(l+1)z}{\prod_{j\leq l+1}(1+jz/E)}=\frac{E}{\prod_{j\leq l}(1+jz/E)}.
\]
So $\Ph_1=\cone^{-1}(\Thp_i+z\theta\log v)\Srow=E\cone^{-1}(\theta v/v)\Srow^{[1]}$. In the same way, the $l$-th term of $\cone^{-1}(\theta v/v)\Srow^{[1]}$ is $\cone^{-1}\theta v\,r_lz^lv^{l-1}/\prod_{j\leq l}(1+jz/E)$; on it the operator acts as multiplication by $(E+lz)\theta v/v+z\theta\log(\cone^{-1}\theta v)$, and
\[
\frac{E+lz}{\prod_{j\leq l}(1+jz/E)}=\frac{E}{\prod_{j\leq l-1}(1+jz/E)},
\]
with the convention of \eqref{eq:set-S} at $l=0$. Together with $\Ph_2=\xi\cone^2[(\Thp_i+z\theta\log v)\Ph_1-\dtwo\Ph_1]$ this gives the second formula.
\end{proof}

%-------------------------------------------------------------------------------
\subsection{Consequences of the closed form}
\label{sec:row-conseq}

\begin{cor}[Laurent form; carrying identities to the conifold]\genusstatus{AI-proposed, unchecked}
\label{cor:row-carry}
\begin{enumerate}
\item At every colour, with $\cone=\ccf$, each $z^n$-coefficient of $\Ph_a^{(i)}$ is a Laurent polynomial in $v_i$ multiplied by $z$-free functions analytic at the conifold (namely $\cone^{\pm1}$, $\Fc$, $\xi$, $\dtwo$). The same holds in the large-radius frame, since the factor $\hb$ depends on $z$ only, and for $\PB_a^{(i)}=\Delta_i^{-1/2}\Ph_a^{(i)}$ up to the constant $\sqrt{3\epsilon_1\epsilon_2E}^{\,-1}$. Note that $\Delta_0^{\pm1/2}$ is proportional to $v_0^{\mp1}$, so it is a factor of this Laurent type and it is not analytic. So each such coefficient is a meromorphic germ at $y=-1/27$, and an identity between such series that holds on a punctured neighbourhood holds as an identity of Laurent expansions in $\tc$.
\item At colour $0$ only, and in frame B only: since $v_0=O(\tc)$, each $[\tc^k]$-coefficient of $\rho^{\nu}\Ph^{(0)}$ ($\nu\in\bbZ$) involves only finitely many $v$-orders. So $[z^n]$ and $[\tc^k]$ commute, and each $[\tc^k]$-coefficient is a rational function of $z$. This fails at colours $1,2$, where $v_j(0)\neq0$ and $[\tc^k]$ of $\Srow(v_j,z)$ is only a formal power series in $z$; and it fails in the large-radius frame, since $\hb$ is not rational.
\end{enumerate}
\end{cor}

\begin{proof}
By \cref{lem:row-closed12}, $\Ph_1$ and $\Ph_2$ are built from $\Srow^{[1]}$, $\Srow^{[2]}$, with $z$-coefficients polynomial in $v$, from $\theta v=1/E-v/3$, and from $1/v$, $1/v^2$. Further $\theta\log(\cone^{-1}\theta v)=-\Fc-\frac13$, since $\theta(\theta v)=-\theta v/3$. The functions $\cone^{\pm1},\Fc,\xi,\dtwo$ are analytic at $\xi=0$ at $\cone=\ccf$ by \cref{lem:set-coord}(a), and $v_i$ is analytic at $\xi=0$ on the branch of \cref{sec:set-colours}. This gives (i). For (ii), by the same formulas every term $r_l(vz)^l/\prod$ carries a power of $v$ growing with $l$, the denominators being rational in $z$; as $v_0=\tc\cdot(\text{$z$-free unit})$ and $\rho$ is a $z$-free unit, only finitely many $l$ enter at each order in $\tc$.
\end{proof}

\begin{cor}[Consequences of the closed form]\genusstatus{AI-proposed, unchecked}
\label{cor:row-closed}
Write $v=v_0$ where colour $0$ is meant.
\begin{description}
\item[(F2)] $\Srow(0,z)=E/(E+z)$: only $l=0$ survives at $v=0$.
\item[(L1)] The $l$-th summand of $\PB_0^{(0)}=v\Srow/\sqrt{3\epsilon_1\epsilon_2E}$ carries $v^{l+1}$ and the proper fraction $z^l/\prod_{j\leq l+1}(1+jz/E)$. At $z=-\psi$, its partial fractions lie in $\mathrm{span}\{\lf_1,\ldots,\lf_{l+1}\}$. As $v=\tc\cdot(\text{$z$-free unit})$, it follows that
\[
[\tc^k]\PB_0^{(0)}\in\mathrm{span}\{\lf_1,\ldots,\lf_k\},\qquad[\tc^k](\rho\Srow)\in\mathrm{span}\{\lf_1,\ldots,\lf_{k+1}\}
\]
(one more in the second, because of the division by $\tc$). These are two different objects. Neither has a constant or a polynomial part.
\item[(Form)] $[z^k](v\,\hb\,\Srow)$ is a polynomial in $v$ of degree $\leq k+1$; equivalently it has the shape $\mathsf{N}_k(Y)/(E+3Y)^{k+1}$ with $\deg\mathsf{N}_k\leq k$.
\item[(R)] For $i=1,2$, every $z$-coefficient of the rows $\RLR_a^{(i)}$, $\PB_a^{(i)}$ ($a=0,1,2$) with $\cone=\ccf$ is analytic at $y=-1/27$. They are polynomials in $\cone^{\pm1}$, $\Fc$, $\dtwo$ with coefficients rational in $t$, regular where $t(1+3\omega^it)\neq0$ and $1-\omega^i\neq0$; further, $\ccf$ and $\Fc$ are analytic with $\ccf(0)=-1$, and $\xi\dtwo=E(\cone^{-2}-1)/2$.
\item[(Z-infinity)] Coefficientwise in $v=v_0$,
\begin{equation}
\label{eq:row-zinf}
\lim_{z\to\infty}z\,v\,\Srow(v,z)=\Sinf=-\tc,\qquad\Sinf\coloneqq\sum_{l\geq0}\frac{r_l\,(Ev)^{l+1}}{(l+1)!}.
\end{equation}
Equivalently $\lim_{z\to\infty}z\rho\Srow=E$.
\end{description}
\end{cor}

\begin{proof}
(F2), (L1) and (Form) are read off from \eqref{eq:set-S}, with $v$ a polynomial variable; for (Form), recall that $\hb$ depends on $z$ only and that $[z^j]\Srow$ has degree $\leq j$ in $v$. (R) follows from \cref{def:set-rows}, since at $t=-1/3$ one has $1+3\omega^it=1-\omega^i\neq0$, together with \cref{lem:set-coord}(a).

For (Z-infinity): coefficientwise in $v$, $f\coloneqq v\Srow=\Sinf/z+O(z^{-2})$. Expand $\Pconj=\sum_{m=0}^{3}z^m\Pconj_m$ at $\alpha=\alpha^*$. Since $\Pconj$ is $\LPF$ with $z\theta$ replaced by $z\theta+Y$, one has $\Pconj_0=\Lhat(Y)=0$, so $\Pconj=z\Pconj_1+z^2\Pconj_2+z^3\Pconj_3$. The top part is
\[
\Pconj_3=\xi\theta^3+27y\theta^2+6y\theta=\Ltwo\circ\theta.
\]
The $V$-adic form of \cref{thm:row-pf} gives $\Pconj f=0$ coefficientwise in $v$, as identities of rational functions in $z$. The operators $\Pconj_1,\Pconj_2$ have poles at $V=0$ (for example $\Pconj_1=-(W/Y)(\theta+\frac12\theta\log W)$, see \cref{lem:row-wkb}(i); only the fact that their coefficients are rational in $V$ is used, and it can be read off from $Y=E/V-E/3$, $y=(V-3)^{-3}$), so everything is read coefficientwise in Laurent series in $V$: each $V$-coefficient of $\Pconj f$ is a finite combination of $V$-coefficients of $f$, with coefficients rational in $z$, and each $V$-coefficient of $f$ is $O(1/z)$. So only $z^3\Pconj_3$ contributes in the limit, and taking $z^{-2}\Pconj f$ as $z\to\infty$ gives $\Ltwo(\theta\Sinf)=0$. Equivalently: $\mathsf{e}_n\to V^n/n!$, $\mathsf{s}\to\dd/\dd V$, $\mathsf{N}\to V\dd/\dd V$, $\mathsf{t}\to3\theta$, so that $\Pconj/z^3\to\theta^3-(\dd/\dd V)^3=\Ltwo\circ\theta$.

The exponents of $\Ltwo$ at $\xi=0$ are $(0,0)$, so $\theta\Sinf=\kappa\,\ccf=\kappa\,\theta\tc$ for a constant $\kappa$, and $\Sinf=\kappa\tc$ because $\Sinf$ vanishes at $v=0$. The $v^1$-coefficients ($Ev$ against $-Ev$, \cref{lem:set-coord}(b)) give $\kappa=-1$. The equivalent form follows from $\rho=-Ev/\tc$.
\end{proof}

The identity \eqref{eq:row-zinf} is exactly the place where the normalisation of $\tc$ enters: it identifies $\Sinf(v_0)=-\tc$ with $\tc$ normalised by $\ccf=\theta\tc$ and $\ccf(0)=-1$.

\begin{prop}[Symplecticity]\genusstatus{AI-proposed, unchecked}
\label{prop:row-sympl}
Let $\epsilon_1,\epsilon_2$ be complex numbers with $\epsilon_1\epsilon_2E\neq0$, or indeterminates. For any non-vanishing section $\cone$ with $\Fc=\theta\log\cone$, and all $i,j$,
\[
\sum_{a,b}\RLR_a^{(i)}(z)\,\eta^{ab}\,\RLR_b^{(j)}(-z)=\delta_{ij}.
\]
The same holds for $\PB_a^{(i)}$, since $\hb(z)\hb(-z)=1$, and so
\[
\big\langle\Ph^{(i)}(z),\Ph^{(j)}(-z)\big\rangle=\delta_{ij}\Delta_i.
\]
By step \textup{(S0)} of the proof, the left side depends only on row $0$, and each of its $z$-coefficients is rational in $t$. So the identity holds at each $z$-order as an identity of rational functions of $t$, and its consequences at the conifold follow by expanding in $\tc$, with no continuation or branch argument.
\end{prop}

The proof uses only \cref{thm:row-pf} and algebra on the closed form. It uses no genericity of $\epsilon_1,\epsilon_2$, so no specialisation argument is needed.

\begin{proof}
Let $\Bbbk$ be the constant field ($\bbC$, or $\bbQ(\omega,\epsilon_1,\epsilon_2)$), and work over $\Bbbk(t)$, with $y=t^3$ and $\theta=(t/3)\dd/\dd t$, with coefficients series in $z$. Step (S0) takes place in a differential extension of $\Bbbk(t)$ containing $\cone^{\pm1}$, $\Fc$ and a square root $\kappa_\circ$ of $3\epsilon_1\epsilon_2E$; steps (S1)--(S6) take place in $\Bbbk(t)\fps{z}$ only. Recall $\Thp_i=Y_i+z\theta$. At colour $i$ put
\[
\mathrm{a}\coloneqq v_i\Srow(v_i,z),\qquad\mathrm{b}\coloneqq\Thp_i\mathrm{a},\qquad\mathrm{c}\coloneqq\xi\,\Thp_i\mathrm{b}.
\]
Let
\[
\mathscr{B}(\mathrm{u},\mathrm{w})\coloneqq-\frac{E^3}{27}\mathrm{u}_0\mathrm{w}_0+E\mathrm{u}_1\mathrm{w}_1+\mathrm{u}_1\mathrm{w}_2+\mathrm{u}_2\mathrm{w}_1,
\]
so that $\eta^{ab}=-3\epsilon_1\epsilon_2\mathscr{B}$ by \eqref{eq:set-etaTF}. Put $\widetilde\Pi_{ij}(z)\coloneqq\mathscr{B}\big((\mathrm{a},\mathrm{b},\mathrm{c})^{(i)}(z),(\mathrm{a},\mathrm{b},\mathrm{c})^{(j)}(-z)\big)$.

\emph{(S0) Drop-out, for any $\cone$.} Put $\mathsf{Q}_1\coloneqq\Thp_i\PB_0$. Using $\theta\cone^{-1}=-\cone^{-1}\Fc$ and $\xi\dtwo=E(\cone^{-2}-1)/2$, \cref{def:set-rows} gives $\PB_1=\cone^{-1}\mathsf{Q}_1$ and
\[
\PB_2=\cone\big[\mathsf{T}_1-z\xi\Fc\mathsf{Q}_1\big]-\cone^{-1}\frac E2\mathsf{Q}_1,\qquad\mathsf{T}_1\coloneqq\xi\Thp_i\mathsf{Q}_1+\frac E2\mathsf{Q}_1.
\]
In the pairing, the $\cone^{-2}$ terms from $\eta^{11}$ and $\eta^{12}$ cancel, using $\eta^{11}=E\eta^{12}$, and so do the $\Fc$ terms. This requires the same $\cone$ at both colours, which holds since $\cone$ depends on $y$ only. There remains
\[
\eta^{00}\PB_0^{(i)}\PB_0^{(j)\prime}-3\epsilon_1\epsilon_2\big(\mathsf{Q}_1^{(i)}\mathsf{T}_1^{(j)\prime}+\mathsf{T}_1^{(i)}\mathsf{Q}_1^{(j)\prime}\big),
\]
with $'$ meaning $z\mapsto-z$. We have $\PB_0=\mathrm{a}/\kappa_\circ$, $\mathsf{Q}_1=\mathrm{b}/\kappa_\circ$ and $\mathsf{T}_1=(\mathrm{c}+\frac E2\mathrm{b})/\kappa_\circ$, and the factors $\hb$ contribute $\hb(z)\hb(-z)=1$. So the frame-B pairing equals $-3\epsilon_1\epsilon_2\widetilde\Pi_{ij}/\kappa_\circ^2=-\widetilde\Pi_{ij}/E$, for every $\cone$. It remains to prove that $\widetilde\Pi_{ij}=-E\delta_{ij}$.

\emph{(S1) First-order system.} Conjugating \eqref{eq:set-pf} at $\alpha=\alpha^*$ gives the stripped form
\[
\xi\Thp^3+(E+27yz)\Thp^2+\Big(\frac{E^2}{3}+6yz^2\Big)\Thp+\frac{E^3}{27},
\]
which annihilates $\mathrm{a}$ by \cref{thm:row-pf}. Since $\Thp(\xi\mathsf{f})=\xi\Thp\mathsf{f}+27yz\mathsf{f}$, the terms $27yz$ cancel, and
\begin{equation}
\label{eq:row-S1}
\Thp\mathrm{a}=\mathrm{b},\qquad\Thp\mathrm{b}=\frac{\mathrm{c}}{\xi},\qquad\Thp\mathrm{c}=-\frac E\xi\mathrm{c}-\Big(\frac{E^2}{3}+6yz^2\Big)\mathrm{b}-\frac{E^3}{27}\mathrm{a}.
\end{equation}
Write this as $\Thp(\mathrm{a},\mathrm{b},\mathrm{c})=\mathcal{D}(\mathrm{a},\mathrm{b},\mathrm{c})$. The matrix $\mathcal{D}$ depends only on $y$ and $z^2$, not on the colour and not on the sign of $z$.

\emph{(S2) Self-adjointness.} For any vectors $\mathrm{u},\mathrm{w}$,
\[
\mathscr{B}(\mathcal{D}\mathrm{u},\mathrm{w})=-\frac{E^3}{27}(\mathrm{u}_1\mathrm{w}_0+\mathrm{u}_0\mathrm{w}_1)+\frac{\mathrm{u}_2\mathrm{w}_2}{\xi}-\Big(\frac{E^2}{3}+6yz^2\Big)\mathrm{u}_1\mathrm{w}_1,
\]
because the two terms $E\mathrm{u}_2\mathrm{w}_1/\xi$ cancel. This is symmetric in $\mathrm{u},\mathrm{w}$, so $\mathscr{B}(\mathcal{D}\mathrm{u},\mathrm{w})=\mathscr{B}(\mathrm{u},\mathcal{D}\mathrm{w})$.

\emph{(S3) Differential equation.} Let $\mathrm{u}\coloneqq(\mathrm{a},\mathrm{b},\mathrm{c})^{(i)}(z)$ and $\mathrm{w}\coloneqq(\mathrm{a},\mathrm{b},\mathrm{c})^{(j)}(-z)$. By \eqref{eq:row-S1} at $z$ and at $-z$: $z\theta\mathrm{u}=\mathcal{D}\mathrm{u}-Y_i\mathrm{u}$ and $z\theta\mathrm{w}=Y_j\mathrm{w}-\mathcal{D}\mathrm{w}$. Since $\mathscr{B}$ has constant coefficients, (S2) gives
\[
(Y_i-Y_j+z\theta)\widetilde\Pi_{ij}=0.
\]

\emph{(S4) The case $i\neq j$.} $Y_i-Y_j$ is a non-zero element of $\Bbbk(t)$ and $\widetilde\Pi_{ij}$ lies in $\Bbbk(t)\fps{z}$, so induction on the $z$-order gives $\widetilde\Pi_{ij}=0$.

\emph{(S5) The case $i=j$.} $\theta\widetilde\Pi_{ii}=0$, so each $z$-coefficient is constant in $t$. Colour $i$ is the image of colour $0$ under the field automorphism $t\mapsto\omega^it$, which fixes $y$, $\xi$ and $\theta$. Hence $\widetilde\Pi_{ii}=\widetilde\Pi_{00}\eqqcolon\kappa(z)$ for all $i$.

\emph{(S6) The constant.} Each $z$-coefficient of $\mathrm{a}$ lies in $v\Bbbk[v]$. Next,
\[
\mathrm{b}=\Big(1-\frac{Ev}3\Big)\Srow+z\Big(\frac1E-\frac v3\Big)\partial_v(v\Srow),
\]
using \eqref{eq:row-thetav}, so each $z$-coefficient of $\mathrm{b}$ lies in $\Bbbk[v]$. And $\mathrm{c}=(\xi Y)\mathrm{b}+\xi z\theta\mathrm{b}$, with $\xi Y=-(E/3)\xi+Et(1-3t+9t^2)$. All are regular at $t=-1/3$, that is $v_0=0$. There $\mathrm{a}=0$; $\mathrm{b}=\Srow(0,z)(1+z/E)=1$, since $\Srow(0,z)=E/(E+z)$; $\xi=0$; $\xi Y=-E$; so $\mathrm{c}=-E$. Thus
\[
\kappa(z)=E\cdot1\cdot1+1\cdot(-E)+(-E)\cdot1=-E,
\]
exactly in $z$.
\end{proof}

\emph{The edge in $\delta$-form.} By \cref{prop:row-sympl}, $\cP_{ij}(-w,w)=-\langle\Ph^{(i)}(w),\Ph^{(j)}(-w)\rangle=-\delta_{ij}\Delta_i$, which does not depend on $w$. So \eqref{eq:set-edge} becomes
\begin{equation}
\label{eq:row-edgedelta}
\Edgeh_{ij}(\psi,w)=\frac{\delta_{ij}\Delta_i-\big\langle\Ph^{(i)}(-\psi),\Ph^{(j)}(-w)\big\rangle}{\psi+w},
\end{equation}
which is the Givental edge in $\Delta$-normalised form. Since $\eta$ is symmetric, this gives the \emph{edge symmetry} $\Edgeh_{ij}(\psi,w)=\Edgeh_{ji}(w,\psi)$ for all $i,j$.

%-------------------------------------------------------------------------------
\subsection{The relations between the generators}
\label{sec:row-modrel}

\begin{prop}\genusstatus{AI-proposed, unchecked}
\label{prop:row-modrel}
Let $A,D,\Ey,X$ be defined by \eqref{eq:set-I} and \eqref{eq:set-gens}. Then, as elements of $\bbQ\fps{y}$,
\[
A^2D=X,\qquad2\Ey=D-X.
\]
Moreover,
\begin{equation}
\label{eq:row-DEX}
\frac{2}{e_1(\alpha)}\theta I_1^{[1]}-\theta I_1^{[2]}+A\log\frac{y}{1+27y}=0,
\end{equation}
and for every $d\geq1$,
\begin{equation}
\label{eq:row-conv}
\sum_{k=0}^{d-1}\frac{1}{d-k}\,\frac{(\frac13)_k(\frac23)_k}{(k!)^2}=\frac{a_d\,(3H_{3d}-H_d)}{27^d}.
\end{equation}
So the two relations are identities of meromorphic germs on $|y|<1/27$ (holomorphic where $A\neq0$; no use is made of the location of zeros of $A$) and along any continuation.
\end{prop}
\draftnote{BS labels: lem:modrel, eq:DEXrel, eq:limitconv. BS writes the right side of the last identity with digamma values; it equals the form above, by $\psi^{(0)}(3d)-\psi^{(0)}(d+1)=H_{3d-1}-H_d$ and $\frac1d+3H_{3d-1}=3H_{3d}$.}

\begin{proof}
\emph{Step 0: reading off \eqref{eq:set-I}.} With $s=1/z$, the $d$-th summand ($d\geq1$) of $I/(zy^{H/z})$ is
\[
y^d\frac{\prod_{m=0}^{3d-1}(-m-3Hs)}{\prod_j\prod_{m=1}^{d}(m+(H+\alpha_j)s)}=3H(-y)^d\frac{(3d-1)!}{(d!)^3}\,s\,\Big[1+s\big(3HH_{3d-1}-(3H+e_1(\alpha))H_d\big)+O(s^2)\Big].
\]
Multiplying by $zy^{H/z}$ and using $3d(3d-1)!=(3d)!$ gives
\[
A=\sum_da_d(-y)^d,\qquad G\coloneqq\frac{\theta I_1^{[1]}}{e_1(\alpha)}=-\sum_da_dH_d(-y)^d,\qquad P_I\coloneqq\theta I_1^{[2]}=A\log y+\hI,
\]
with
\[
\hI\coloneqq\sum_d3a_d(H_{3d}-H_d)(-y)^d.
\]
For $\hI$: the term from $I_0^{[1]}$ contributes $(a_d/d)(-y)^d$, and $1/d+3H_{3d-1}=3H_{3d}$. So $\Ey=\theta(G/A)$, $D=\theta(P_I/A)$, and $\Ey$ does not depend on $\alpha$. No reduction modulo the cubic relation occurs at this order. This agrees with \eqref{eq:set-Ian}.

\emph{Step 1: $\Ltwo$ on series.} For $\varphi=\sum\varphi_d(-y)^d$,
\[
[(-y)^d]\Ltwo\varphi=d^2\varphi_d-3(3d-1)(3d-2)\varphi_{d-1}.
\]
Since $d^2a_d=3(3d-1)(3d-2)a_{d-1}$:
\begin{enumerate}
\item $\Ltwo A=0$;
\item if $\varphi_d=a_d\mathrm{g}_d$, then $[(-y)^d]\Ltwo\varphi=d^2a_d(\mathrm{g}_d-\mathrm{g}_{d-1})$;
\item $\Ltwo$ is injective on $y\bbQ\fps{y}$, because the lowest term $d\geq1$ has coefficient $d^2\varphi_d\neq0$.
\end{enumerate}
We also use the product rule, for a function $f$: $\Ltwo(Af)=\xi(2\theta A\,\theta f+A\theta^2f)+27yA\theta f$.

\emph{Step 2: the key identity.} We claim
\begin{equation}
\tag{K}
\sum_da_d(3H_{3d}-H_d)(-y)^d=-A\log(1+27y).
\end{equation}
Both sides lie in $y\bbQ\fps{y}$. For the left side, $\mathrm{g}_d=3H_{3d}-H_d$ gives $\mathrm{g}_d-\mathrm{g}_{d-1}=3/(3d-1)+3/(3d-2)=9(2d-1)/((3d-1)(3d-2))$, so by (ii) the $(-y)^d$-coefficient of $\Ltwo$ of the left side is $27(2d-1)a_{d-1}$. For the right side, $\theta\log\xi=27yX$ and $\theta^2\log\xi=27yX^2$, so the product rule gives $\Ltwo(A\log\xi)=54y\theta A+27yA$, whose $(-y)^d$-coefficient is $-27(2d-1)a_{d-1}$. Hence $\Ltwo$ of the right side has the same coefficient as $\Ltwo$ of the left side. Apply (iii). Comparing coefficients in (K), after dividing by $27^d$, is exactly \eqref{eq:row-conv}, since $a_k/27^k=(\frac13)_k(\frac23)_k/(k!)^2$.

\emph{Step 3: \eqref{eq:row-DEX}.} We have $2G-\hI=-\sum a_d(3H_{3d}-H_d)(-y)^d=A\log\xi$ by (K). So $2G-P_I+A\log(y/\xi)=0$, which is \eqref{eq:row-DEX}.

\emph{Step 4: the second relation.} Divide by $A$: $G/A=\frac12(P_I/A-\log(y/\xi))$. Since $\theta\log(y/\xi)=1-27yX=X$, applying $\theta$ gives $2\Ey=D-X$.

\emph{Step 5: the first relation.} We have $\Ltwo P_I=0$. Indeed, by (ii) with $\mathrm{g}_d=3(H_{3d}-H_d)$, $[(-y)^d]\Ltwo\hI=27(2d-1)a_{d-1}-2da_d$. By the product rule, $\Ltwo(A\log y)=2\xi\theta A+27yA$, with coefficient $2da_d-27(2d-1)a_{d-1}$. The sum is $0$. Put $\mathsf{Wr}\coloneqq A\theta P_I-P_I\theta A=A^2D$. Then $\theta\mathsf{Wr}=-27yX\,\mathsf{Wr}$, so $\theta(\xi\mathsf{Wr})=0$. Now $\mathsf{Wr}=A^2+A\theta\hI-\hI\theta A$ is a log-free series with $\mathsf{Wr}(0)=1$. So $\xi\mathsf{Wr}=1$.
\end{proof}

\begin{rmk}[Dictionary at $\cone=A$]
\label{rmk:row-dict}
Recall $\dtwo=E(\cone^{-2}-1)/(2\xi)$, with $E$ the weight. At $\cone=A$ we have $\xi\cone^2=A^2/X=1/D$ by the first relation of \cref{prop:row-modrel}, and
\[
\dtwo=\frac E2\Big(\frac{X}{A^2}-X\Big)=\frac E2(D-X)=e_1(\alpha^*)\Ey
\]
by both relations. So \cref{def:set-rows} at $\cone=A$ has the recursion coefficients of the row recursion of \cite{BS}: the rows there satisfy $\Rm_{1j}=(z\theta\Rm_{0j}+Y_j\Rm_{0j})/A$ and $\Rm_{2j}=(z\theta\Rm_{1j}+(Y_j-e_1(\alpha)\Ey)\Rm_{1j})/D$. The identification of row $0$ is a separate statement, \cref{lem:spe-row0,lem:spe-red}. \draftnote{BS label: eq:RHb.}
\end{rmk}
